\section{NDN 背景与命名职责}
\subsection{命名获取、转发状态与对象身份}
NDN 通过 Interest 和 Data 按名称获取数据~\cite{ndn,ndnpacket}。应用发出描述所需 Data 的 Interest，生产者或中间缓存可提供匹配的 Data。转发器根据转发信息选择 Interest 的方向，并利用待满足 Interest 的状态交付返回的 Data。因此，消费者无需指定数据当前的存储位置，就可以复用网络中的命名 Data。可用性仍取决于转发、连接及可达副本；名称并不保证交付。

本文中的名称承担几种相关职责。服务名称标识所请求的操作；消息名称及经过验证的字段标识调用交换；对象名称标识图像、模型分片或输出张量。名称也参与约束可接受签名密钥的信任规则~\cite{trustschema}。这些用途需要一致，但不能混同：匹配服务前缀并不证明服务权限，有效生产者签名也不证明对象就是当前执行被分配的输入。

例如，一张图像可以采用以下示意名称：
\begin{center}\texttt{/org/uav7/view/task42/view2/version1}\end{center}
它属于采集无人机的生产者命名空间。应用定义任务、视角及版本分量的含义，并验证签名者。处理 Provider 可以引用该图像，而不必将它改名为自己的观测。这个例子解释命名原则，并非强制 NDNSF wire format。即使 Content 已加密，名称和流量模式仍可能可见，因此不能仅因载荷受保护就把敏感标识放入公开名称。

转发器的 Content Store 提供机会性缓存；repository 为应用提供显式管理的对象保留位置。缓存命中或摘要检查成功都不等于完整的依赖准入判断。消费者仍需核验生产者权限、目标对象身份及当前调用或协作关联。同样，Data 的 FreshnessPeriod 是指导缓存新鲜度处理的 NDN 元数据；它不是服务期限，不能替代消息到期、当前权限检查或应用防重放状态。

Sync 帮助参与者发现命名数据集的变化~\cite{li2018brief,moll2022sok}，应用随后获取并验证相应 Data。因此，NDNSF 的“发布”指使命名对象可被获取，并通过配置的传播机制通告其可用性，不应理解为未经 Interest 请求就向任意转发器主动推送 Data。Sync 恢复可在时间允许时补充缺失通告或触发获取，但不能保证任务在截止时间之后完成，也不能克服无限期断连。

这也明确了本文 data-driven execution 的含义。经过验证的 Request、ACK、Selection、结果或声明的中间对象，为本地状态迁移提供依据。计时器和本地调度策略仍然存在：超时可以关闭收集或终止操作，规划算法可以决定发布何种分配。影响其他参与者的决定通过命名、可验证的消息传达；仅在本地运行回调不能通知远程参与者。这不意味着全部运行状态都必须成为网络 Data，也不意味着数据到达就足以授权执行。

\subsection{从 NDN 原语到服务决策}
Forwarding Information Base（FIB）提供名称前缀转发信息，Pending Interest Table（PIT）记录待满足的获取请求，Content Store（CS）保存缓存 Data~\cite{yi2013case}。兼容的待满足 Interest 可以共享返回的 Data。这种网络层聚合不同于选择多个服务执行者。PIT 条目的生命周期也不同于服务操作：条目到期不会自动取消应用 handler 或删除 User 的请求状态。

\begin{table}[htbp]\centering\small
\caption{NDN 基础机制及其上层仍需承担的服务决策。}
\label{tab:ndn-runtime-map}
\begin{tabularx}{\textwidth}{p{.23\textwidth}XX}
\toprule 基础机制 & NDNSF 用法 & 额外服务职责\\\midrule
Interest/Data 与名称 & 获取消息、结果及对象分段 & 定义操作、请求身份及完成条件\\\midrule
FIB、PIT、CS & 转发和复用命名 Data & 选择工作、跟踪生命周期、检查对象版本\\\midrule
签名与 trust schema & 验证签名权限和 Data 完整性 & 检查服务权限及本次调用的分配\\\midrule
数据集同步（泛化 Sync；NDN-SVS 实现） & 发现新发布的消息名称 & 验证消息并处理 ACK、Selection 和超时状态\\\midrule
内容加密 & 限制受保护 Content 的读取者 & 定义接收者、分发密钥、应用权限变更\\\bottomrule
\end{tabularx}
\end{table}

软件栈中，本地 NDN 转发器负责包转发，ndn-cxx 提供应用使用的包、签名、验证和获取接口~\cite{ndn-cxx}；NDN State Vector Sync（NDN-SVS）提供具体的数据集同步实现~\cite{ndnsvs-github}。NDNSF 将这些设施组织为服务交互。服务策略不要求转发器解释 GPU 负载、飞行就绪状态或协作角色；这些含义由框架和应用契约定义。

\section{远程过程调用（RPC）、机器人中间件与消息机制}
\subsection{RPC 接口与微服务集成}
远程过程调用（remote procedure call，RPC）框架为跨进程或跨主机操作提供类型化接口。gRPC 是支持单次调用和流式调用的一种 RPC 实现~\cite{grpc}。ROS 等机器人中间件为机器人组件提供可复用的通信抽象~\cite{ros,ros2}，面向消息的中间件则支持生产者与消费者解耦~\cite{messagequeue}。这些系统都是替代应用自行编写网络代码的相关方案。NDNSF 不声称端点发现、异步交互或分布式应用组合在已有中间件中不存在。

端点解析能够找到目标，但不一定能证明该目标有权并且当前愿意执行某次调用。中间件或应用可以另外加入这些策略；本文将这些策略如何表达、以及调用方能够获得哪些证据，作为比较的一部分。

对比应关注如何表达相同应用需求。RPC 实验需要说明端点如何解析、连接如何管理、如何重试以及如何检测故障；NDN 实验需要说明名称如何发现、配置的同步和重传如何工作、Data 如何验证以及如何处理 Content Store。类型化调用并不天然比命名对象交换更不安全，安全性取决于配套认证、授权和加密。同样，网络缓存只有在工作负载和操作语义允许复用对象时才有价值；复用图像并不意味着具有物理副作用的调用可以安全重复执行。

NDNSF 需要回答的是：如何通过名称和可验证对象保持服务决策与应用数据的关联。评价必须在明确的工作负载下检验这种关联，而不能把网络抽象的差异直接等同于更优。下面讨论的命名计算与命令授权系统提供了更直接的协议先例。

\subsection{机器人组件与异步消息}
机器人软件将相机、遥测、控制和感知分成组件，常把服务调用与持续更新的发布—订阅 topic 结合使用~\cite{ros,ros2}。NDNSF 的服务容器设计采用组件行为与通信支持分离的思路，并不替代飞控，也不暗示机器人中间件不能实现授权。比较重点是进入系统的操作如何指向正确车辆，以及保存后的观测如何继续关联其生产者与操作。

消息 broker 为日志、遥测和工作流事件提供异步交付与缓冲~\cite{messagequeue}。相关问题包括数据保留位置、发布/消费权限，以及交付确认和任务完成的区别。稳定部署可以适合使用 broker。NDN 命名获取提供不同的放置和验证方式，但延迟获取仍需要明确的保留策略。两类消息抽象本身都不决定哪个 Provider 愿意执行，也不定义重复物理命令是否安全。

\subsection{微服务与边缘计算}
CFEC 利用移动性和资源信息研究车辆 fog、edge 与 cloud 间的计算卸载~\cite{cfec}。MIA-NDN 根据输入特征研究面向微服务的 Interest 聚合和 PIT 生命周期管理~\cite{miandn}。SECaaS 研究在 multi-access edge 部署访问控制实施和匿名化等安全功能~\cite{secaas}。它们将比较范围扩展到函数调用之外的资源放置、转发状态和可部署安全功能。

不能把这些系统的不同目标直接改写为 NDNSF 的独有能力。放置可行性不等于调用权限，感知参数的聚合不决定协作中的指定前驱，部署访问控制服务也不同于授权调用该服务。反过来，NDNSF 的应用层检查不替代这些系统的资源管理或转发机制。声称性能更优需要共同工作负载和匹配需求；它们的公开结果不能直接作为本文实验的可互换基线。
